% Électrolyse d'une solution de bromure de cuivre(II) (exercice 2, chapitre 10)
\begin{tikzpicture}[x=1cm,y=1cm]
\pic at (0,0) {becher={largeur=4.6,hauteur=2.5,niveau=2.0,liquide=pcBleu!35}};
\fill[orange!55!brown!40,opacity=0.7] (1.1,0.28) ellipse (0.55 and 0.22);
\fill[orange!55!brown!30,opacity=0.6] (1.1,0.8) ellipse (0.4 and 0.45);
\filldraw[pcAppareil] (-1.22,0.3) rectangle (-0.98,2.9);
\filldraw[pcAppareil] (0.98,0.3) rectangle (1.22,2.9);
\fill[pcCuivre] (-1.26,0.3) rectangle (-0.94,1.35);
\draw[pcFil] (-1.1,2.9) -- (-1.1,3.5) -- (-0.45,3.5);
\draw[pcFil] (1.1,2.9) -- (1.1,3.5) -- (0.45,3.5);
\pcAppareilRond{(0,3.5)}{G}
\draw[pcCourant] (1.75,3.5) -- (1.1,3.5);
\node[pcLegende,anchor=west,text=otRed!80!black] at (1.8,3.5) {$I$};
\pcNote[east]{(-1.26,1.0)}{(-2.2,1.2)}{électrode 1 : dépôt rougeâtre}
\pcNote[west]{(1.45,0.7)}{(2.4,0.9)}{électrode 2 : coloration orangée}
\node[pcLegende,anchor=north] at (0,-0.15) {solution de bromure de cuivre(II)};
\end{tikzpicture}
\par\smallskip\legende{Les observations après quelques minutes d'électrolyse ; les bornes du générateur ne sont pas indiquées.}
